% Part: history
% Chapter: biographies 
% Section: alan-turing 

\documentclass[../../../include/open-logic-section]{subfiles}

\begin{document}

\olfileid{his}{bio}{tur} 

\olsection{اېلن ټيورينګ}

اېلن ټيورينګ د 1912 د جون پر 23مه
د لندن په مېډا وېل کښې زېږېدلے و۔
هغه د نظري کمپيوټرپوهنې پلار ګڼل
کېږي۔ په طبيعي علومو او رياضياتو
کښې ئې شوق له کم عمره پيل شو۔
خو په ماشومتوب کښې ئې ښوونځيو
دې شوقونو ته مناسب پام نۀ کولو،
ځکه هلته په ادبياتو او کلاسيکو
زده کړو ټينګار کېده۔ په نتيجه
کښې په ښوونځي کښې کمزورے و او
ډېرو ښوونکو به ترټلو۔

\olphoto{turing-alan}{اېلن ټيورينګ}

ټيورينګ د کېمبرېج په کينګز کالج کښې
د لومړۍ پوهنتوني درجې دپاره
رياضيات ولوستل۔ په 1936 کښې ئې
هغه څه جوړ کړل چې اوس د ټيورينګ
ماشين بلل کېږي، څو د محاسبه
کېدونکې تابع مفهوم په کره ډول
تعريف کړي او د پرېکړې د مسئلې
ناپرېکړه کېدنه ثابته کړي۔ الونزو
چرچ تر هغه مخکښې دې نتيجې ته
رسېدلے و، او په خپل لامبډا حساب
ئې ثابته کړې وه۔ د ټيورينګ مقاله
بيا هم د چرچ نتيجې ته په حوالې
سره خپره شوه۔ چرچ ټيورينګ
پرينسټن ته راوبالۀ؛ هغه له
1936--1938 پورې هلته و او د
چرچ تر لارښوونې لاندې ئې
دوکتورا واخيسته۔

په منطق کښې له شوق سره سره، په طبيعي
علومو کښې د ټيورينګ پخوانے شوق هم
غښتلے پاتې شو۔ د دويمې نړيوالې
جګړې پر وخت د بلېچلي پارک په
برتانوي د رمزشننې څانګه کښې د
کار پر مهال ئې عملي مهارتونه
وکارول شول۔ ټيورينګ د جرمني د
سمندري پوځي اړيکو د رمز، يعنې
اېنيګما، په ماتولو کښې مرکزي
شخص و۔ په احصائيه او رمزپوهنه
کښې د هغه مهارت، او د الېکټروني
ماشينونو پېژندل، ډلې ته دا
وړتيا ورکړه چې د «بامب» په نوم
د رمز پرانيستونکے ماشين جوړ
کړي او رمز مات کړي۔ د هغه
نظرونو د نړۍ د لومړي پروګرام
منونکي الېکټروني کمپيوټر،
کولوسس، په جوړېدو کښې هم مرسته
وکړه؛ دا هم په بلېچلي پارک کښې
د جرمني د لورېنز د رمز د ماتولو
دپاره کارېده۔ \emph{سرچينه‌يي
يادونه OLSTH-045:} پورته متن د
بامب د جوړېدو بيان له الېکټروني
ماشينونو سره ګډوي۔ د کمپيوټرپوهنې
ملي موزيم د بامب جوړښت برقي او
ميخانيکي ښيي؛ کولوسس بيا
الېکټروني و۔ د دواړو ماشينونو
جوړښت او د اېنيګما او لورېنز
کارونه بايد جلا وساتل شي۔

ټيورينګ همجنس‌پال و۔ سره له دې، په
1942 کښې ئې د بلېچلي پارک خپلې
همکارې جون کلارک ته د وادۀ وړانديز
وکړ، خو وروسته ئې کوژده ماته
کړه او هغې ته ئې څرګنده کړه چې
همجنس‌پال دے۔ په ژوند کښې ئې
څو معشوقان لرل، که څه هم هغه
وخت همجنسي اعمال په برتانيه کښې
جرمي سرغړونې وې۔ په 1952 کښې
ئې د هغه وخت د معشوق يو ملګري
کور غلا کړ، او د پوليسو د راپور
د ورکولو پر وخت ټيورينګ ومنله
چې همجنسي اړيکه لري؛ هغه ګڼله
چې حکومت دا اعمال قانوني کوي۔
دا سمه نۀ وه، او د سختې
بې‌حيايۍ تور پرې ولګېد۔ زندان
ته د تلو پر ځای ټيورينګ هغه
هورموني درملنه غوره کړه چې
جنسي ميل ئې کمولو۔ ټيورينګ د
1954 د جون پر 8مه مړ وموندل
شو؛ علت ئې د ساينايډ زيات مقدار
و، او غالباً ځان‌وژنه وه۔ په
2013 کښې ملکې الېزابت~II ورته
شاهي بخښنه ورکړه۔ \emph{سرچينه‌يي
يادونه OLSTH-046:} د کوژدې کال
او د څرګندونې ترتيب د منجمد متن
پورته بيان دے۔ د ټيورينګ د
ژوندليکوال اېنډريو هاجز آرشيف
د وادۀ وړانديز د 1941 په پيل
کښې، د همجنس‌پالۍ څرګندونه ورسره
سمدستي، او د کوژدې ماتېدل بيا
د هماغه کال په اګست کښې ښيي۔

\begin{reading}
د اېلن ټيورينګ د يوه هراړخيز
ژوندليک دپاره \citet{Hodges2014}
وګورئ۔ د هغه ژوند او کار د
\emph{د رمز ماتول} ډرامې الهام
شو؛ دا په 1996 کښې د ټېلېويژن
دپاره جوړه شوه، او ډېريک جکوبي
پکښې د ټيورينګ رول ولوباوه۔
\emph{د تقليد لوبه} هم، چې
بېنېډېکټ کمبربېچ او کېرا نايتلي
پکښې لوبېږي او د اکاډمۍ جايزې
ته نومول شوے فلم دے، د اېلن
ټيورينګ پر ژوند او په بلېچلي
پارک کښې پر وخت په آزاد ډول
ولاړ دے \citep{Imitation2014}۔

\citet{Radiolab2012} د ټيورينګ د
ژوند او کار په هکله څو غږيزې
خپرونې لري۔ د بي بي سي د
هورايزن مستند فلم \emph{د ډاکټر
ټيورينګ عجيب ژوند او مرګ} په
انټرنېټ کښې کتل کېدلے شي
\citep{Sykes1992}۔ \citep{Theelen2012}
د ليګو د يوه کارکوونکي ټيورينګ
ماشين لنډه ويډيو ده؛ دا په
2012 کښې د هغه د زېږون د
سلمې کليزې د درناوي دپاره
جوړ شوے و۔

د ټيورينګ د ماشينونو او د پرېکړې
د مسئلې په هکله د هغه اصلي
مقاله \citet{Turing1937} ده۔
\end{reading}

\end{document}
